% JAC1964-20261121. /root/programme_audit; GPT-6 Astra Pro.
% Accepted input: numerical=homological on all abelian varieties,032Cor8.4.
% Additional established input: Jannsen semisimplicity.

\subsection*{JAC-P16. Lefschetz inverses for varieties of geometric abelian type}

\begin{theorem}[JAC-P16]\label{jac:P16}
Let $k$ be any field and let $X/k$ be smooth projective of dimension
$d$. Suppose that $h(X_{\bar k})$ belongs to the thick rigid
category of rational Chow motives generated by abelian varieties
and Artin motives. Accept the equality of numerical and
$\ell$-adic homological equivalence on every abelian variety over
an algebraically closed field, for every invertible prime $\ell$,
as stated in 032, Corollary~8.4, printed p.~50.

Then numerical and homological equivalence agree on $X^a$ in
every codimension, for every $a\ge1$. In particular this assertion
holds on $X\times X$, not only for cycles on $X$.
The Kunneth components of the diagonal of $X$ are algebraic.
For every ample class $h$ defined over $k$ and every $0\le i\le d$,
the inverse of
\[
 L_h^{d-i}:H^i(X_{\bar k},\mathbf Q_\ell)
       \longrightarrow
 H^{2d-i}(X_{\bar k},\mathbf Q_\ell(d-i))                 \tag{JAC.P16.1}
\]
is induced by a rational algebraic correspondence defined over $k$.
One can choose these correspondences to work for every
$\ell\ne\operatorname{char}k$. The normalized lowering operator
$\Lambda_X$ is algebraic as well; thus $B(X)$ holds.

If $k$ is algebraically closed, JAC-P02 consequently proves all
three restriction/fixed-trace assertions of cote~44, Th.~6, for
every such $X$ and a genuine Lefschetz pencil on $X$, in every
characteristic. This includes its deleted critical finite-kernel
clause. No Hodge-standard assertion for $X$'s own intersection
form is included in this theorem.
\end{theorem}

\begin{proof}
We first work over $\bar k$ and fix one invertible prime $\ell$.
Let $\mathcal S$ be the category whose objects are finite direct
sums of Tate twists of motives of abelian varieties. Products and
duals remain of this form after adding direct summands. Artin
motives are sums of units over $\bar k$, and the unit is already
an abelian-type summand. Therefore every Chow motive of abelian
type is a direct summand of an object of $\mathcal S$.

\emph{Numerical and homological ideals agree on the whole category.}
For $S,T\in\mathcal S$, a morphism is a matrix of algebraic
cycle correspondences on products of abelian varieties, in the
codimensions determined by the Tate twists. Those products are
themselves abelian varieties. The accepted equality in 032
therefore identifies the numerically zero and homologically zero
matrices. Equivalently, the categorical numerical ideal
\[
 \mathcal N(S,T)
  =\{f:\operatorname{tr}(gf)=0\text{ for every }g:T\to S\}
                                                               \tag{JAC.P16.2}
\]
is exactly the kernel of realization. This is not an assumption
about just endomorphisms: choosing $g$ with one matrix component
tests each component of $f$ against all complementary cycles.
Under Poincare duality, the trace is their intersection number.
Jannsen's trace identity is Lemma~1, printed p.~448, of
\emph{Motives, numerical equivalence, and semi-simplicity},
Invent.\ Math.\ 107 (1992), 447--452.

Here is the direct-summand step explicitly. Suppose $M$ is a
summand of $S$ and $N$ a summand of $T$, with inclusions $u_M,u_N$
and retractions $v_M,v_N$. For $f:M\to N$ extend it by zero to
\[
                 F=u_N f v_M:S\longrightarrow T.
\]
For every $G:T\to S$, cyclicity of the trace gives
\[
 \operatorname{tr}_S(GF)
     =\operatorname{tr}_M(v_MG u_N f).                    \tag{JAC.P16.3}
\]
Every $g:N\to M$ occurs as $v_MG u_N$, by taking
$G=u_Mg v_N$. Thus $f$ is numerically zero if and only if $F$ is.
The ambient equality of ideals makes this equivalent to $H(F)=0$,
which is in turn equivalent to $H(f)=0$, since the inclusions and
retractions split after realization. This proves the equality of
the two ideals on all abelian-type summands and all their Hom
spaces, including twists.

For a variety $W$ of abelian type, codimension-$r$ cycles are
morphisms from the appropriate Tate motive to $h(W)$; testing
their categorical numerical ideal is exactly testing intersection
with complementary cycles on $W$. Hence the previous paragraph
proves numerical=homological on $W$ in every codimension.
Apply this to $W=X^a$, whose motive is of abelian type for every
$a$. In particular, all codimensions on $X\times X$ have the
required equality. This supplies the full correspondence-level
statement used below.

\emph{Algebraic Kunneth projectors.}
Choose a Chow splitting $h(X)\xrightarrow{u}S\xrightarrow{v}h(X)$
with $vu=1$. Abelian varieties have algebraic Kunneth projectors,
with their functorial grading; a direct sum with Tate twists has
the induced degree projectors $\varepsilon_i$. The correspondences
\[
                         \pi_i=v\varepsilon_i u          \tag{JAC.P16.4}
\]
act as the degree-$i$ projections on $H^*(X)$. They are therefore
orthogonal idempotents modulo homological equivalence, and also
modulo numerical equivalence by the equality just proved. We
need only algebraicity of their cohomological classes, not a claim
that the particular Chow representatives in (JAC.P16.4) are
orthogonal idempotents before passing to that equivalence relation.

\emph{Semisimplicity gives the inverse, without a positivity argument.}
The pseudo-abelian category of homological abelian-type motives
is now fully faithfully identified with the full numerical
subcategory generated by the same objects: the equality of ideals
was proved on all Hom spaces, and idempotent completion preserves
that equality. Jannsen's Theorem~1, printed p.~448, makes the
numerical category semisimple abelian. Its full subcategory closed
under direct sums and summands is also semisimple abelian.

Using (JAC.P16.4), (JAC.P16.1) is a morphism between the two
appropriate degree motives, with its indicated Tate twist.
Hard Lefschetz makes its realization an isomorphism. In a
semisimple category every short exact sequence splits, so any
additive realization is exact. Its faithfulness forces both kernel
and cokernel of this morphism to vanish. The morphism is therefore
invertible in the numerical category.

The inverse is a rational correspondence class. Lift it to a
cycle on $X\times X$ with the prescribed degree supports. The
two inverse identities initially hold numerically; the proved
equality on $X\times X$ makes them hold homologically. Thus the
lift induces precisely the hard-Lefschetz inverse, rather than
just some map with a numerically inverse class.
The construction is independent of the selected prime at the
numerical-motive level. The same Kunneth classes work in every
realization, and 032 gives the equality of ideals for every
invertible prime. The lifted inverse therefore works for all of
them.

For clarity, the normalized lowering operator can be reconstructed
by finite algebraic compositions. For $r\le d$ let
$I_r=(L^{d-r})^{-1}:H^{2d-r}(X)\to H^r(X)$, with the degree
supports understood. The primitive projector on $H^r(X)$ is
\[
 e_r=1-L I_{r-2}L^{d-r+1}\quad(r\ge2),
 \qquad e_0=e_1=1,                                     \tag{JAC.P16.5}
\]
where $I_{r-2}=(L^{d-r+2})^{-1}$.
It kills $LH^{r-2}$ and is the identity on
$\ker L^{d-r+1}$. For $0\le a\le d-r$, the projector onto
$L^aP^r\subset H^{r+2a}(X)$ is
\[
             L^a e_r I_r L^{d-r-a}\pi_{r+2a}.           \tag{JAC.P16.6}
\]
On that string, the lowering map for $a\ge1$ is obtained by
replacing the first $L^a$ in (JAC.P16.6) by $L^{a-1}$; it is
zero for $a=0$. Summing these finitely many algebraic
correspondences gives the normalized $\Lambda_X$ used in the
archive and JAC-P02.

\emph{Descent to the original field.}
Each correspondence constructed over $\bar k$ is defined over a
finite extension $L/k$. Its desired cohomological action is
canonical: a Kunneth projection, a hard-Lefschetz inverse for
the $k$-defined class $h$, or the resulting normalized lowering
operator. It is therefore Galois invariant. If
$\pi:(X\times X)_L\to X\times X$, replace a representing
correspondence $\gamma$ by
\[
                         [L:k]^{-1}\pi_*\gamma.          \tag{JAC.P16.7}
\]
After geometric base change, this is the average of the conjugate
correspondences, with their inseparable multiplicities. Each has
the same canonical action. The number of embeddings is the
separable degree and each multiplicity is the inseparable degree,
so their sum is $[L:k]$; (JAC.P16.7) has exactly the required
action. This deals with imperfect $k$ without pretending every
finite extension is Galois or separable.

The numerical/homological equality also descends. A numerically
trivial $k$-cycle remains numerically trivial after any finite
extension, by the projection formula against the norm of an
arbitrary complementary cycle. Every complementary cycle over
$\bar k$ is defined over a finite extension. Hence a numerically
trivial cycle is geometrically numerically trivial and therefore
homologically zero by the geometric result. The reverse implication
always follows from intersection compatibility. This applies
to $X^a$ in every codimension.

We have proved $B(X)$. Over an algebraically closed field, insert
the normalized algebraic $\Lambda_X$ into JAC-P02. Its proof
already establishes the below-middle isogeny, the critical image
and fixed trace, the critical finite kernel, and above-middle
surjectivity. No change to its cycle-spread or imperfect-pencil-field
argument is needed.
\end{proof}

\paragraph{Why the Hodge-standard form is still a separate question.}
The proof establishes algebraic inverses and nondegeneracy of the
pairings on the relevant numerical summands. It does not identify
the geometric primitive form
\[
       (x,y)\longmapsto(-1)^r\deg_X(xyh^{d-2r})           \tag{JAC.P16.8}
\]
with a positive form transported from an abelian variety. Even
if a positive form $Q_0$ is supplied on an abstract summand,
another nondegenerate symmetric form has the shape
$Q(x,y)=Q_0(ax,y)$ with $a$ invertible and $Q_0$-self-adjoint.
These two properties alone do not make its eigenvalues positive;
$a=-1$ already demonstrates that logical gap. Here one must prove
positivity for the particular comparison operator determined by
(JAC.P16.8). The Chow-motive presentation and Jannsen
semisimplicity supply no such assertion. This is not a
counterexample to Hodge standard for abelian-type varieties;
it records exactly why that stronger conclusion has not been
claimed from this proof.

\paragraph{Primary source and archive scope.}
Jannsen's original six-page article was directly read at
\url{https://epub.uni-regensburg.de/26642/1/jannsen10.pdf}:
Theorem~1 and the trace identity are on printed p.~448;
the proof of semisimplicity continues on p.~449.
The functorial abelian Kunneth decomposition is recalled as
Deninger--Murre's theorem in Ancona--Enright-Ward--Huber,
Theorem~4.2.2, printed p.~824, in the source checked for JAC-P14.
The accepted numerical/homological input is explicitly the first
assertion of 032, Corollary~8.4, printed p.~50.
This provides a precise extension of cote~44, scan~4, Th.~6,
and of the inverse-correspondence discussion in cote~16,
scan~49. Neither a universal regular representative nor
Hodge-standard positivity for every such $X$ is inferred.
